\section{El regreso del relato del Pre-training}

La sección anterior expuso el hilo conductor de todo el episodio; esta sección aborda la primera postura no consensuada. Después de 2024, la discusión en la industria se desplazó en gran medida hacia el Post-training, el RL, los reasoning model de la serie o, las aplicaciones de Agent y la conversión en producto, por lo que se volvió común el juicio de que «los rendimientos del Pre-training se están desacelerando». Aquí Guangmi (广密) sostiene lo contrario: el Pre-training no ha terminado y sigue siendo la etapa más esencial para elevar el límite superior de nuevas capacidades.

\begin{figure}[H]
\centering
\includegraphics[width=0.92\textwidth]{pretrain-posttrain-split.png}
\caption{Pre-training y Post-training: el preentrenamiento determina el límite superior intrínseco; el posentrenamiento y el RL se parecen más a dar forma a las capacidades. Diagrama conceptual propio, elaborado a partir de la conversación entre 00:04:22 y 00:11:56.}
\end{figure}

\begin{knowledgebox}{Lectura de la figura: el límite superior y la forma de las capacidades no son el mismo problema}
A la izquierda, el Pre-training se encarga de comprimir a gran escala el conocimiento del mundo, el código, los patrones de razonamiento y las capacidades latentes; a la derecha, el Post-training/RL se encarga de alinear esas capacidades con tareas, preferencias, herramientas y límites de seguridad. El juicio de Guangmi es que el posentrenamiento puede hacer que el modelo «use mejor las capacidades que ya tiene», pero difícilmente puede crear de la nada capacidades de las que carece el modelo base.
\end{knowledgebox}

\subsection{Por qué el Pre-training sigue siendo una postura no consensuada}

Esta sección desglosa de dónde surge lo «no consensuado». Del lado del consenso, mucha gente observó que después de GPT-4 el ritmo de lanzamiento de modelos se hizo más lento, aumentaron las discusiones sobre el muro de datos y los benchmark de RL/reasoning subieron de puntuación con rapidez, de modo que concluyó que los beneficios del preentrenamiento del modelo base están por agotarse. Del lado no consensuado, Guangmi considera que la desaceleración temporal de un líder no equivale al fin de un paradigma; si el equipo de Pre-training de OpenAI atraviesa inestabilidad, desde fuera podría confundirse la situación organizacional de una empresa con el límite técnico de toda la línea de desarrollo.

\begin{importantbox}{Guangmi enfatiza: el líder no es la línea de desarrollo misma}
Cuando una empresa líder se desacelera en cierta línea, la industria tiende a interpretarlo como que «esta línea técnica se terminó». Pero también puede tratarse solo de cambios en la capacidad organizacional, la rotación de talento, la asignación de recursos y las prioridades estratégicas. Para juzgar si una línea técnica ha terminado, hay que observar el desempeño conjunto de varias empresas, varias arquitecturas y la siguiente generación de modelos.
\end{importantbox}

\noindent\begin{tabularx}{\textwidth}{p{0.20\textwidth}p{0.34\textwidth}X}
\toprule
Nivel de análisis & Dónde es fácil equivocarse & Lectura más prudente \\
\midrule
Avance de una empresa & El Pre-training de OpenAI se desacelera en cierta etapa & Puede ser un problema organizacional y estratégico; no equivale a que toda la industria se desacelere. \\
Benchmark & Los reasoning model suben de puntuación con rapidez & Las puntuaciones suben rápido, pero eso no necesariamente eleva todo el límite superior de capacidades. \\
Muro de datos & Hay cada vez menos texto de internet de alta calidad & Los datos sintéticos, los entornos de código y la retroalimentación de herramientas pueden cambiar la oferta de datos. \\
Interés por las aplicaciones & Los Agent y las aplicaciones atraen más atención & El auge de las aplicaciones sigue dependiendo de que la inteligencia del modelo base siga empujando hacia arriba. \\
\bottomrule
\end{tabularx}

\subsection{Qué puede hacer y qué no puede hacer el Post-training y el RL}

Esta sección aclara el papel del Post-training/RL. Guangmi no niega el posentrenamiento; se opone a tratarlo como sustituto de las capacidades del modelo base. El posentrenamiento puede hacer que el modelo siga mejor las instrucciones, sea más seguro, maneje mejor herramientas específicas y razone mejor en tareas de matemáticas y código; el RL puede, mediante señales de recompensa, activar patrones de razonamiento que el modelo ya posee. Pero si el conocimiento, las representaciones y los priores de acción del base model son insuficientes, el posentrenamiento se parece a «un alumno de primaria que resuelve ejercicios en serie»: la calificación sube a corto plazo, pero el techo a largo plazo es limitado.

\begin{knowledgebox}{Términos clave: Pre-training, Post-training, RL}
\noindent\begin{tabularx}{\textwidth}{p{0.18\textwidth}p{0.36\textwidth}X}
\toprule
Término & Mecanismo & Papel en este episodio \\
\midrule
Pre-training & Entrenar el base model con datos a gran escala para aprender representaciones generales y capacidad de predicción & Determina el límite superior intrínseco del modelo; es el hilo conductor que Guangmi vuelve a enfatizar. \\
Post-training & Dar forma al comportamiento del modelo con datos de instrucciones, preferencias, seguridad y tareas & Mejora la usabilidad, pero sobre todo activa y organiza capacidades existentes. \\
RL & Optimizar el comportamiento mediante señales de recompensa; se usa con frecuencia en reasoning, código y tareas con herramientas & Puede reforzar los patrones de exploración y resolución de problemas, pero depende de las capacidades del modelo base. \\
Synthetic Data & Datos de entrenamiento generados por un modelo o un sistema & Podría aliviar el muro de datos, en especial las trayectorias de CoT/herramientas de alto valor. \\
\bottomrule
\end{tabularx}
\end{knowledgebox}

\begin{warningbox}{Error común: confundir «rendir mejor en exámenes» con «ser más inteligente»}
Los reasoning model avanzan rápido en matemáticas, código y benchmark, pero eso no significa que ya hayan adquirido todas las capacidades nuevas. Hay que distinguir entre «aprovechar mejor capacidades existentes» y «que las representaciones subyacentes realmente mejoren».
\end{warningbox}

\subsection{Las dificultades de los datos sintéticos y de los marcos de entrenamiento}

Esta sección explica con más detalle por qué el Pre-training todavía puede continuar. Guangmi menciona que los datos de CoT de alto valor, los datos generados por RL y el sampling dentro de entornos podrían volver a incorporarse a la etapa de preentrenamiento y aliviar el data bottleneck. Pero no se trata simplemente de volver a meter el texto generado en el conjunto de entrenamiento: exige que el marco de entrenamiento integre de forma más estrecha inference, sampling, reward, filtering y pre-training. Dicho de otro modo, el preentrenamiento del futuro podría dejar de ser un entrenamiento sobre un corpus estático y convertirse en una fábrica de datos con exploración, generación y retroalimentación.

\begin{importantbox}{Experiencia práctica: el muro de datos no se resuelve solo con «más páginas web»}
Si en el futuro los datos de alto valor provienen de la exploración del modelo, la ejecución de código, las llamadas a herramientas y la retroalimentación del entorno, la ingeniería de datos pasará de «limpiar texto de internet» a «ejecutar entornos, muestrear trayectorias, evaluar resultados y reincorporarlos al entrenamiento». Esa es también la razón por la que los Agent y el Online Learning regresan al hilo central del entrenamiento de modelos.
\end{importantbox}

\subsection{Resumen de la sección}

El regreso del relato del Pre-training no significa que el posentrenamiento carezca de importancia, sino que el posentrenamiento no puede sustituir el límite superior del modelo base. La pregunta decisiva del primer trimestre de 2025 es qué empresas todavía pueden seguir ampliando las capacidades del base model y cuáles solo están empaquetando capacidades existentes en productos y benchmark de mejor apariencia.
